% =====================================================================
\phantomsection
\addcontentsline{toc}{part}{Languages of the Web}
\BookChapter[JavaScript and TypeScript]{JavaScript and TypeScript: The Languages of the Web}
\label{ch:javascript}

\begin{researchquestion}
How do you turn a calculation into something a colleague can use in a web browser?
\end{researchquestion}
\OrientationNote
% =====================================================================

\begin{factsheet}{JavaScript}
First appeared & 1995; standardised as ECMAScript (ES), yearly editions since ES2015\\
Designed by & Brendan Eich, Netscape\\
Paradigms & imperative, functional, object-oriented (prototype-based, with classes since ES2015), event-driven\\
Typing & dynamic, weak; TypeScript adds static types\\
Execution & JIT-compiled by the browser's engine (V8, SpiderMonkey, JavaScriptCore) or by Node.js, Deno, Bun\\
File extensions & \file{.js}, \file{.mjs}; TypeScript \file{.ts}\\
Used for & every interactive web page; web applications; dashboards and data visualisation;
  servers (Node.js); desktop apps (Electron); scripting in some engineering tools\\
Famous software & Google Maps, VS Code (TypeScript), D3.js, Plotly.js, Three.js, React,
  Jupyter's web interface\\
\end{factsheet}

JavaScript is the only programming language that runs natively in every web browser,
which makes it the most widely \emph{deployed} language in the world. Every time you use
an interactive web page --- an online map, a booking form, an online simulation, your
university's learning-management system --- JavaScript is running on your computer. For
engineers, it is the language of web dashboards, interactive visualisations, online
calculators and the user interfaces of many modern tools.

\LearningOutcomes{Run JavaScript in a browser console and with Node.js; avoid the common surprises of loose typing; process data with array methods and JSON; build a small interactive calculator in a single web page; explain event-driven and asynchronous code; and say what TypeScript adds.}

\section{History and philosophy}

\begin{history}
In May 1995 Netscape, maker of the dominant web browser, asked Brendan Eich to create a
simple scripting language for web pages. He produced the first version in about ten days.
It was called Mocha, then LiveScript, and finally JavaScript --- a marketing decision to
ride on the popularity of Java, with which it has little in common. To prevent each
browser from developing its own dialect, the language was standardised as
\emph{ECMAScript} in 1997. In 2008 Google released the V8 engine, whose JIT compiler made
JavaScript fast enough for large applications, and in 2009 Ryan Dahl built
\textbf{Node.js} on V8 so that JavaScript could also run outside the browser, on servers.
\end{history}

JavaScript was designed to be forgiving: a web page should keep working even if a script
has a small mistake. This led to features --- automatic type conversions, a special value
\code{undefined} for missing things, functions that accept any number of arguments --- that
make the language easy to start with but can hide errors in calculations. Modern
JavaScript (since ES2015, ``ES6'') added \code{let}/\code{const}, arrow functions, classes,
modules and many other improvements, and \textbf{TypeScript} adds a static type system on
top (Section~\ref{sec:typescript}).

\section{Getting started}

You already have a JavaScript environment: open any web browser, press \code{F12} (or
\code{Ctrl+Shift+I}, \code{Cmd+Option+I} on a Mac) and choose the ``Console'' tab. Type
\code{2 + 3} and press Enter. For running scripts from files, install
\textbf{Node.js} from \url{https://nodejs.org}:
\codefile{JavaScript}{javascript/hello.js}{hello world}
\begin{lstlisting}[style=shell]
$ node hello.js
\end{lstlisting}
\outfile{javascript/hello.out}
Strings in backticks are \term{template literals}: \code{\$\{\ldots\}} inserts the value of
an expression, like an f-string in \py.

\section{The language in brief}

\subsection*{Declare with \texttt{const} and \texttt{let}; every number is a double}

Declare variables with \code{const} (cannot be reassigned) or \code{let} (can be
reassigned). Prefer \code{const}; avoid the old keyword \code{var}, which has confusing
scoping rules.
\begin{lstlisting}[language=JavaScript]
const g = 9.81;                 // number
let n = 0;                      // number (there is no separate integer type)
const fluid = "water";          // string
const ok = true;                // boolean
let x;                          // undefined (declared, no value yet)
const nothing = null;           // deliberately empty
const big = 2n ** 64n;          // BigInt: arbitrary-precision integer
typeof g;                       // "number"
\end{lstlisting}
All ordinary numbers are IEEE double-precision floats, so \code{7 / 2} is \code{3.5} and
integers are exact only up to $2^{53}\approx 9\times10^{15}$. Mathematical functions live in
the \code{Math} object (\code{Math.sqrt}, \code{Math.exp}, \code{Math.sin},
\code{Math.PI}, \ldots) and powers can be written \code{x ** 2}.

\subsection*{Control flow follows C, and functions are values}

Control flow follows C: \code{if}/\code{else}, \code{for}, \code{while}, \code{switch},
\code{break}, \code{continue}, braces and (optional, but recommended) semicolons. Use
\code{for (const item of array)} to loop over the values of an array --- the counterpart of
\py's \code{for item in list}.

Functions can be declared in two ways:
\begin{lstlisting}[language=JavaScript]
function terminalVelocity(m, c, g = 9.81) {      // function declaration
  return Math.sqrt(m * g / c);
}
const area = (d) => Math.PI * d ** 2 / 4;        // arrow function
\end{lstlisting}
Functions are values: they can be stored in variables, passed to other functions and
returned. This makes functional programming natural in JavaScript.

\subsection*{Always compare with three equals signs}

JavaScript has two equality operators. \code{==} converts types before comparing
(\code{"5" == 5} is \code{true}); \code{===} does not (\code{"5" === 5} is \code{false}).
Always use \code{===} and \code{!==}. The automatic conversions of the \code{+} operator
are a frequent source of bugs when numbers come from text input fields, which always
produce strings: \code{"3" + 4} is \code{"34"}. Convert explicitly with
\code{Number(text)} or \code{parseFloat(text)}.

\subsection*{Arrays and objects carry the data; JSON carries it between programs}

\codefile{JavaScript}{javascript/arrays.js}{arrays, objects and JSON}
\outfile{javascript/arrays.out}
Array methods such as \code{map}, \code{filter} and \code{reduce} take functions as
arguments and are the idiomatic way to process data. An \term{object literal}
\code{\{ name: "S275", density: 7850 \}} is JavaScript's dictionary. Its text form is
\term{JSON} (JavaScript Object Notation), which has become the universal data-exchange
format of the web and is read by every language in these notes.

The last line of the output shows three classic surprises. \code{"3" + 4} joins strings;
\code{"3" * 4} converts to numbers; and \code{sort()} without arguments sorts numbers
\emph{as strings}, so \code{10} comes before \code{9}. Always pass a comparison function,
\code{sort((a, b) => a - b)}, when sorting numbers.

\section{Case study: the falling skydiver}

\codefile{JavaScript}{javascript/drag.js}{Euler solution of the drag equation}
\outfile{javascript/drag.out}
The numbers are identical to those of all the other languages; only the exponent format of
\code{toExponential} differs slightly. JavaScript has no \code{printf}, so we built two
small helper functions with \code{toFixed} and \code{padStart}. \code{Float64Array} is a
\term{typed array}: a compact block of doubles, much like a NumPy array, which the JIT
compiler can process efficiently.

\section{JavaScript in the browser}

JavaScript's real home is the web page. A page consists of three layers: \textbf{HTML}
describes the content and structure, \textbf{CSS} describes the appearance, and
\textbf{JavaScript} provides the behaviour. Scripts read and change the page through the
\term{Document Object Model} (DOM), and they react to \term{events} such as clicks, key
presses and slider movements.

The following file is a complete interactive engineering tool: an online calculator that
re-runs the drag simulation and redraws the plot every time the user moves a slider. Save
it as \file{drag.html} and open it in any browser --- no installation, no server.
\codefile[language=HTML]{HTML}{javascript/drag.html}{an interactive web calculator}
This \emph{event-driven} style is different from the scripts we have seen so far. The
program does not run from top to bottom and stop; it sets up the page, registers a function
to be called when the slider moves (\code{addEventListener}), and then waits for the user.
Graphical user interfaces in every language (C\# Windows applications, \py's Tkinter or
Qt, MATLAB App Designer) work the same way.

\begin{engnote}[Engineering note: share a calculation as a web page]
A single HTML file with a little JavaScript is often the easiest way to give colleagues or
clients a calculation tool: everyone has a browser, nothing needs to be installed, and the
file can be emailed or put on an intranet. For serious plots, use a library such as
Plotly.js, Chart.js or D3.js instead of drawing on a canvas by hand. And remember that
anyone can read the code of a web page --- and change it in their copy.
\end{engnote}

\subsection*{Asynchronous code keeps the page responsive while it waits}

Web pages must not freeze while waiting for data from a server. JavaScript therefore uses
\term{asynchronous} functions: a request is started, the program continues, and the result
is handled when it arrives. Modern code uses \code{async}/\code{await}:
\begin{lstlisting}[language=JavaScript]
async function loadSensorData(url) {
  const response = await fetch(url);          // wait without blocking the page
  if (!response.ok) throw new Error(`HTTP ${response.status}`);
  const data = await response.json();          // parse the JSON body
  return data.readings.map((r) => r.temperature);
}

loadSensorData("https://example.com/api/plant/42/readings")
  .then((temps) => console.log("Max temperature:", Math.max(...temps)))
  .catch((err) => console.error("Could not load data:", err));
\end{lstlisting}
This pattern is how dashboards display live data from sensors, databases and simulations.

\section{TypeScript}
\label{sec:typescript}

\textbf{TypeScript} (Microsoft, 2012, also designed by Anders Hejlsberg of C\# fame) is
JavaScript with optional static types. A TypeScript compiler checks the types and then
removes them, producing ordinary JavaScript. Most large JavaScript projects --- including
VS Code itself --- are now written in TypeScript, because the type checker catches errors
before the code runs.
\begin{lstlisting}[language=TypeScript]
interface Material {
  name: string;
  density: number;        // kg/m^3
  youngsModulus: number;  // Pa
}

function specificStiffness(m: Material): number {
  return m.youngsModulus / m.density;
}

const al: Material = { name: "Al6061", density: 2700, youngsModulus: 69e9 };
specificStiffness(al);              // OK
specificStiffness({ name: "x" });   // compile-time error: properties missing
const s: number = "25";             // compile-time error: string is not a number
\end{lstlisting}

\section{Beyond the browser}

\begin{itemize}
\item \textbf{Node.js}, \textbf{Deno} and \textbf{Bun} run JavaScript and TypeScript on
  servers and on the command line, with access to files and the network. The package
  registry \textbf{npm} is the largest software repository in the world.
\item \textbf{Electron} packages a web application as a desktop program; VS Code, Slack
  and many other applications are built this way.
\item \textbf{WebAssembly} (Wasm) lets code written in C, C++, Rust, Fortran or even \py
  run in the browser at near-native speed, with JavaScript as the glue. This is how
  JupyterLite and Pyodide run \py --- including NumPy --- inside a web page.
\item JavaScript is the scripting language of some engineering and office tools, such as
  Google Apps Script (Google Sheets), Adobe Acrobat forms and Node-RED, a visual tool for
  wiring together Internet-of-Things devices.
\item Visualisation libraries such as \textbf{D3.js}, \textbf{Plotly.js} and
  \textbf{Three.js} (3D graphics) are the basis of many interactive scientific figures; the
  \py and R versions of Plotly generate JavaScript underneath.
\end{itemize}

\section{Strengths and weaknesses}

\begin{center}
\small
\begin{tabularx}{\linewidth}{@{}>{\raggedright\arraybackslash}X>{\raggedright\arraybackslash}X@{}}
\toprule
\textbf{Strengths} & \textbf{Weaknesses}\\
\midrule
Runs in every browser; nothing to install & Weak typing causes surprising results\\
The only choice for interactive web pages & Only one number type (double); no built-in matrices\\
Fast JIT engines & Small numerical/scientific ecosystem\\
Huge ecosystem (npm) and community & Ecosystem changes very fast; dependency sprawl\\
TypeScript adds safety for large projects & Event-driven and asynchronous style takes practice\\
\bottomrule
\end{tabularx}
\end{center}

\section{Further reading}
\begin{itemize}
\item MDN Web Docs (Mozilla), \url{https://developer.mozilla.org/en-US/docs/Web/JavaScript} ---
  the best reference and tutorial.
\item M.~Haverbeke, \emph{Eloquent JavaScript}, 4th ed., No Starch Press, 2024 (free online).
\item The TypeScript Handbook, \url{https://www.typescriptlang.org/docs/}.
\end{itemize}

\begin{exercises}
\item Open your browser's console and evaluate \code{0.1 + 0.2}, \code{"5" - 2},
  \code{"5" + 2}, \code{[] + []}, \code{typeof NaN} and \code{[1, 2, 10].sort()}. Explain
  each result.
\item Open \file{drag.html} in a browser. Add a second slider for the mass and show the
  terminal velocity as a dashed horizontal line on the plot.
\item Write a JavaScript function that converts a temperature in Celsius to Fahrenheit and
  kelvin, and use \code{map} to apply it to an array of temperatures.
\item Use \code{reduce} to compute the mean and the standard deviation of an array of
  numbers in a single pass (hint: accumulate $\sum x$ and $\sum x^2$). Is this one-pass
  formula numerically safe? (Recall Section~\ref{sec:floatingpoint}.)
\item Rewrite the \code{Material} interface example in plain JavaScript. Which errors that
  TypeScript catches would now go unnoticed until run time?
\end{exercises}
